\begin{proof}
    The proof follows by induction.
    \begin{itemize}
    \item \textbf{Base Case}.
          Let $g \in \Sigma_{\Core}^0$. There are six cases to consider.
          \begin{enumerate}
              \item Let $g = \scalebox{0.65}{\circimg{tikz-cache/encoder__eqn3_lhs.pdf}}$.
                    Then
                    $
                        \sem{\encoder(g)}_C
                        =
                        \sem*{ \scalebox{0.65}{\circimg{tikz-cache/encoder__eqn3_rhs.pdf}} }_C
                        =
                        e^{i\alpha\pi}
                        =
                        \exp( \alpha \Log( -1 ) )
                        =
                        \sem{g}_H
                    $.
              \item Let $g = \scalebox{0.65}{\circimg{tikz-cache/encoder__eqn6_lhs.pdf}}$.
                    Then
                    $
                        \sem{\encoder(g)}_C
                        =
                        \sem*{ \scalebox{0.65}{\circimg{tikz-cache/encoder__eqn6_lhs.pdf}} }_C
                        =
                        \tfrac{1}{\sqrt{2}}
                        \begin{bsmallmatrix}
                            1 & 1 \\
                            1 & -1
                        \end{bsmallmatrix}
                        =
                        \sem*{g}_H
                    $.
              \item Let $g = \scalebox{0.65}{\circimg{tikz-cache/encoder__eqn1_lhs.pdf}}$.
                    Recall that $\Log( Z ) = i\pi \ket{1}\bra{1}$.
                    Then the following equation holds.
                    \begin{equation*}
                        \sem{\encoder(g)}_C
                        =
                        \sem*{ \scalebox{0.65}{\circimg{tikz-cache/encoder__eqn1_rhs.pdf}} }_C
                        =
                        \ket{0} \bra{0} \otimes \mathbb{I}_1 + \ket{1} \bra{1} \otimes e^{i\alpha\pi}
                        =
                        \exp( i\alpha\pi \ket{1} \bra{1} )
                        =
                        \sem{g}_H
                    \end{equation*}
              \item Let $g = \scalebox{0.65}{\circimg{tikz-cache/encoder__eqn4_lhs.pdf}}$.
                    Recall that $HZH = X$.
                    Then the following equation holds
                    \begin{align*}
                        \sem{\encoder(g)}_C
                        =
                        \sem*{ \scalebox{0.65}{\circimg{tikz-cache/encoder__eqn4_rhs.pdf}} }_C
                        &=
                        H
                        \exp( \alpha \Log( Z ) )
                        H
                        &
                        ( \text{by case (3)} )
                        \\
                        &=
                        \exp( \alpha H\Log( Z )H )
                        &
                        ( \text{by~\cref{Eqn:UnitaryConj}} )
                        \\
                        &=
                        \exp( \alpha \Log( HZH ) )
                        &
                        ( \text{by~\cref{Eqn:UnitaryConj}} )
                        \\
                        &=
                        \exp( \alpha \Log( X ) )
                        =
                        \sem{g}_H
              \end{align*}
              \item Let $g = \scalebox{0.65}{\circimg{tikz-cache/encoder__eqn2_lhs.pdf}}$.
                    Recall that $XZX = -Z$.
                    Then the following equation holds.
                    \begin{align*}
                        \sem*{\encoder(g)}_C
                        &=
                        \sem*{ \scalebox{0.65}{\circimg{tikz-cache/encoder__eqn2_rhs.pdf}} }_C
                        \\
                        &=
                        \exp( \Log(X) ) \exp( \alpha \Log(Z) ) \exp( \Log(X) )
                        &
                        ( \text{by cases (3) and (4)} )
                        \\
                        &=
                        X \exp( \alpha \Log(Z) ) X
                        \\
                        &=
                        \exp( \alpha X \Log(Z) X )
                        &
                        ( \text{by~\cref{Eqn:UnitaryConj}} )
                        \\
                        &=
                        \exp( \alpha \Log( XZX ) )
                        &
                        ( \text{by~\cref{Eqn:UnitaryConj}} )
                        \\
                        &=
                        \exp( \alpha \Log( -Z ) )
                        =
                        \sem{g}_H
                    \end{align*}
              \item Let $g = \scalebox{0.65}{\circimg{tikz-cache/encoder__eqn5_lhs.pdf}}$.
                    Recall that $XZX = -Z$.
                    Then the following equation holds.
                    \begin{align*}
                        \sem*{\encoder(g)}_C
                        &=
                        \sem*{ \scalebox{0.65}{\circimg{tikz-cache/encoder__eqn5_rhs.pdf}} }_C
                        \\
                        &=
                        \exp( \Log(Z) ) \exp( \alpha \Log(X) ) \exp( \Log(Z) )
                        &
                        ( \text{by cases (3) and (4)} )
                        \\
                        &=
                        Z \exp( \alpha \Log(X) ) Z
                        \\
                        &=
                        \exp( \alpha Z \Log(X) Z )
                        &
                        ( \text{by~\cref{Eqn:UnitaryConj}} )
                        \\
                        &=
                        \exp( \alpha \Log( ZXZ ) )
                        &
                        ( \text{by~\cref{Eqn:UnitaryConj}} )
                        \\
                        &=
                        \exp( \alpha \Log( -X ) )
                        =
                        \sem{g}_H
                    \end{align*}
          \end{enumerate}
          Since $g \in \Sigma_{\Core}^0$ was arbitrary, then $\sem{\encoder(g)}_C = \sem{g}_H$ for all $g \in \Sigma_{\Core}^0$.
    \item \textbf{Inductive Hypothesis}.
          For some $n \in \mathbb{N}$, if $g \in \Sigma_{\Core}^n$, then $\sem{\encoder(g)}_C = \sem{g}_H$.
    \item \textbf{Inductive Step}.
          Assume that the inductive hypothesis holds for some $n \in \mathbb{N}$ and let $g \in \Sigma_{\Core}^{n+1}$.
          If $g \in \Sigma_{\Core}^n$ as well, then $\sem{\encoder(g)}_C = \sem{g}_H$ by the inductive hypothesis.
          Assume instead that $g \in \Sigma_{\Core}^{n+1} \setminus \Sigma_{\Core}^n$.
          Then there exists some $U \in \Mor( \Core[n] )$ such that $g = \scalebox{0.65}{\circimg{tikz-cache/dagger__term1.pdf}} \odot U$.
          Then by the inductive hypothesis, since $\sem{\encoder(-)}_C$ and $\sem{-}_H$ agree on all generators, then $\sem{\encoder(U)}_C = \sem{U}_H$.
          Then by \cref{Thm:GenCtrlNF}, the following holds.
          \begin{equation*}
              \sem{g}_H
              =
              Z \odot \sem{U}_H
              =
              Z \odot \sem{\encoder(U)}_C
              =
              \ket{0} \bra{0} \otimes \mathbb{I}_{2^n} + \ket{1} \bra{1} \otimes \sem{\encoder(U)}_C
              =
              \sem{\encoder(g)}_C
          \end{equation*}
          In either case, $\sem{g}_H = \sem{\encoder(g)}_C$.
          Then the inductive step holds.
    \end{itemize}
    Then by the principle of induction, $\sem{\encoder(g)}_H = \sem{g}_C$ for each $g \in \Sigma_{\Core}$.
\end{proof}
